\section{Simplicial Sets}

\subsection{Simplicial complexes}

Recall that a {\emphcolor simplex} in $\bR^N$ is a subspace,
$$ [v_0,\dots,v_n] = \set{\sum_i\lambda_iv_i}[\sum_i\lambda_i=1,\ 0\leq\lambda_i\in\bR] , $$
where $v_i\in\bR^N$ are {\emphcolor geometrically independent} (if $\sum_i\lambda_iv_i=0$ and $\sum_i\lambda_i=0$ then $\lambda_i=0$).
The {\emphcolor faces} of a simplex $[v_0,\dots,v_n]$ are the simplices generated by a subset of $\set{v_0,\dots,v_n}$.
The {\emphcolor standard $n$-simplex} is the one generated by the standard basis vectors of $\bR^{n+1}$:
$$ \abs{\Delta^n} = [e_0,\dots,e_n] = \set{(x_0,\dots,x_n)\in\bR^{n+1}}[\sum_ix_i=1,\ x_i\geq0] \subseteq \bR^{n+1} . $$

A {\emphcolor simplicial complex} in $\bR^N$ is a collection of simplices $K$ such that the face of any simplex in $K$ is also in $K$, and the intersection of any two
simplices in $K$ is a face of both of them.
Given a simplical complex $K$, we can define its {\emphcolor geometric realization} $\abs K$ to be the union of all the simplices in $K$.
A subset $F\subseteq\abs K$ is closed iff $F\cap\sigma$ is closed in the simplex $\sigma$ for all $\sigma\in K$.

Simplicial complexes are a good starting point, but they are insufficient for our purposes.
A straightforward first generalization is as follows.

\bdefn[title={Abstract simplicial complexes}]

    An {\emphcolor abstract simplicial complex} is a collection of sets $K=\set{X^n}_{n<\omega}$, where $X^k$ is a collection of subsets of $X^0$
    of cardinality $k+1$.
    Furthermore, it must satisfy the condition that any subset of $X^k$ of size $j+1$ belongs to $X^j$.

    Elements of $X^0$ are called the {\emphcolor vertices} of the complex.

\edefn

Alternatively an abstract simplicial complex is a set $K$ of finite subsets of $X^0$ which is closed under taking subsets.
Then $X^n$ is just the collection of sets in $K$ of cardinality $n+1$.

In a simplicial complex of $\bR^N$ we required also that the intersection of simplices be a simplex, but this is because we were dealing with simplices
embedded in an ambient space.

The standard $n$-simplex $\abs{\Delta^n}$ gives rise to the abstract simplicial complex of non-empty subsets of $\set{0,\dots,n}$.

\bdefn[title={Face maps}]

    We will assume that an abstract simplicial complex $K$ comes equipped with a total ordering on its vertices.
    Then for an $n$-simplex $\set{v_0,\dots,v_n}\in K$ we order its vertices and write $[v_0,\dots,v_n]$ to mean $v_0<\cdots<v_n$.

    We further define the {\emphcolor face maps} of the simplicial complex $K$ to be maps $d_0,\dots,d_n\colon X^n\to X^{n-1}$ given by
    $$ d_i[v_0,\dots,v_n] = [v_0,\dots,\hat v_i,\dots,v_n] $$
    where the hat means the vertex was removed.

\edefn

For each $n$ we have maps $d_i$ for $i\leq n$, to differentiate the face maps between simplices of differing dimension we could write $d_i^n$ instead of merely $d_i$,
but this is an overload of information that is not necessary.

There is a fundamental relation between face maps which is immediately clear:
$$ d_id_j = d_{j-1}d_i,\qquad \hbox{for $i<j$} . \eqcnt[faceid] $$
Indeed, both sides map $[v_0,\dots,v_n]$ to $[v_0,\dots,\hat v_i,\dots,\hat v_j,\dots,v_n]$.

\bdefn[title={Simplicial maps}]

    If $K,L$ are simplicial complexes, a {\emphcolor simplicial map} $f\colon K\to L$ is a function from vertices of $K$ to vertices of $L$ such that if
    $\set{v_0,\dots,v_n}\in K$ then $\set{f(v_i)}_i\in L$.
    Note that a simplicial map need not be injective.

\edefn

\subsection{Simplicial sets}

Simplicial sets, like simplicial complexes, can be given a combinatorial definition.
Unfortunately the combinatorial definition is lengthier and instead has a more concise category-theoretic definition.

\bdefn[title={Simplicial sets}]

    The {\emphcolor simplex category} $\Delta$ is the category whose objects are finite ordinals with order-preserving maps between them.
    Explicitly, its objects are $[n]=\set{0,\dots,n}$ and a morphism $f\colon[n]\to[m]$ is a function which preserves order (not necessarily strictly).

    A {\emphcolor simplicial set} is a functor $\Delta^\op\to\Set$.
    Then the category of simplicial sets is the functor category $\sSet=[\Delta^\op,\Set]$.
    Morphisms are natural transformations.

\edefn

Explicitly, a simplicial set $X$ is a collection of sets $X_n=X([n])$ along with functions $f\colon X_n\to X_m$ for every order-preserving $f\colon[m]\to[n]$.
To understand simplicial sets, it is worthwhile to understand the simplex category in more depth.

There are some obvious maps in the simplex category $\Delta$:
\benum
    \item The {\emphcolor face map} $\delta_i\colon[n-1]\to[n]$ (for $i\leq n$) is the unique injection whose image leaves out only $i$.
    Explicitly, $\delta_i(j)=j$ for $j<i$ and $\delta_i(j)=j+1$ for $j\geq i$.
    \item The {\emphcolor degeneracy map} $\sigma_i\colon[n+1]\to[n]$ (for $i\leq n$) is the unique surjection which picks out $i$ twice.
    Explicitly, $\sigma_i(j)=j$ for $j\leq i$ and $\sigma_i(j)=j-1$ for $j>i$.
\eenum
The following identities are then easily verified:
$$ \displaylines{
    \delta_i\circ\delta_j = \delta_{j+1}\circ\delta_i\hbox{ for $i<j$},\cr
    \sigma_j\circ\sigma_i = \sigma_i\circ\sigma_{j+1}\hbox{ for $i\leq j$},\cr
    \sigma_j\circ\delta_i = \cases{\delta_i\circ\sigma_{j-1} & $i<j$\cr \id & $i=j$ or $j+1$\cr \delta_{i-1}\circ\sigma_j & $j+1<i$} .
} $$

Then if $X$ is a simplicial set, we write $X(\delta_i)=d_i\colon X_n\to X_{n-1}$ for the image of the face maps, and $X(\sigma_j)=s_j\colon X_n\to X_{n+1}$ for the image of the degeneracy maps under $X$.
These are still called the face and degeneracy maps.
Note that due to the contravariance of $X$, the face maps ``shrink'' and the degeneracy maps ``grow''.

The identities in the simplex category then dualize to
$$ \vcenter{\displaylines{
    d_i d_j = d_{j-1} d_i\hbox{ for $i<j$},\cr
    s_i s_j = s_{j+1} s_i\hbox{ for $i\leq j$},\cr
    d_i s_j = \cases{s_{j-1} d_i & $i<j$\cr \id & $i=j$ or $j+1$\cr s_j d_{i-1} & $j+1<i$} .
}} \eqcnt[simpid] $$
Note that the first identity is the same as the one in \gotopasteqn{faceid}.

It is not hard to see that every morphism in the simplex category can be written as a composition of face and degeneracy maps.
Thus we can equivalently define a simplicial set to be any collection of sets with face and degeneracy maps satisfying the identities in \gotoleqn.
And a morphism of simplicial sets $X\to Y$ is a sequence of functions $f_n\colon X_n\to Y_n$ which commute with the face and degeneracy maps: $d\circ f=f\circ d$ and $s\circ f=s\circ f$.

If $K$ is a simplicial complex, then it defines a simplicial set $X$ where
$$ \displaylines{
    X_n = \set{[v_0,\dots,v_n]}[\hbox{The $v_i$ form a simplex in $K$ with $v_0\leq\cdots\leq v_n$}],\cr
    d_i[v_0,\dots,v_n] = [v_0,\dots,\hat v_i,\dots,v_n],\cr
    s_i[v_0,\dots,v_n] = [v_0,\dots,s_i,s_i,\dots,v_n].
} $$
Note that the elements of $X_n$ can be {\it degenerate\/} simplices.
For example if $[0,1,2]$ is a simplex of $K$, then $[0,0,1,2,2,2]$ is a simplex of $X$.
Essentially the simplicial set associated with a simplicial complex is obtained by adjoining all of these degenerate simplices.

We can generalize what we mean by a degenerate simplex.

\bdefn

    Let $X$ be a simplicial set.
    A {\emphcolor degenerate simplex} of $X$ is a simplex $x\in X$ which is in the image of some degeneracy map.

\edefn

The prototypical simplicial set is
$$ \Delta^n = \hom_\Delta(\bul,[n]) . $$
We can better characterize $\Delta^n$ as follows.
Notice that an increasing map $f\colon[k]\to[n]$ can be represented as a (degenerate) $k$-simplex $[f(0),f(1),\dots,f(k)]$.
Thus we see that this definition of $\Delta^n$ coincides with the one obtained from taking the simplicial complex defined by $\abs{\Delta^n}$ and adjoining degenerate simplicies.

\subsection{Geometric realization}

Consider the functor $\abs{\Delta^\bul}\colon\Delta\to\Top$ which maps $[n]$ to $\abs{\Delta^n}$, and an order preserving map $\phi\colon[n]\to[m]$ to the continuous map
$\abs{\Delta^n}\to\abs{\Delta^m}$ given by
$$ (x_0,\dots,x_n)\in\abs{\Delta^n} \mapsto (y_0,\dots,y_m)\in\abs{\Delta^m},\qquad y_j = \sum_{\phi(i)=j}x_i . $$
For example, the face maps $\delta_i$ define maps $\abs{\Delta^{n-1}}\to\abs{\Delta^n}$ by
$$ \delta_i(x_0,\dots,x_{n-1}) = (x_0,\dots,0,\dots,x_{n-1}), $$
and the degeneracy maps $\sigma_i\colon\abs{\Delta^n}\to\abs{\Delta^{n-1}}$ give
$$ \sigma_i(x_0,\dots,x_n) = (x_0,\dots,x_i+x_{i+1},\dots,x_n) . $$

This then defines a functor
$$ \hom_\Top(\abs{\Delta^\bul},\bul) \colon \Delta^\op\times\Top \to \Set . $$
In particular for every topological space $X$, we have a simplicial set
$$ \S(X) = \hom_\Top(\abs{\Delta^\bul},X) , $$
this is called the {\emphcolor singular set} of $X$.
The face and degeneracy maps of the singular set are then
$$ d_i(\rho) = \rho\circ\delta_i,\qquad s_i(\rho) = \rho\circ\sigma_i . $$

\bnote

    The singular set of $X$ plays a central role in defining the singular homology of $X$.
    Recall that $X$'s chain complex is defined at each coordinate to be the free Abelian group generated by $\S(X)_n$: $C_n(X)=Z[\S(X)_n]$.
    The differentials are then defined in terms of the face maps: $\partial\colon C_n(X)\to C_{n-1}(X)$ is defined on a singular $n$-simplex $\sigma\in\S(X)_n$ by
    $$ \partial\sigma = \sum_{i=0}^n(-1)^id_i\sigma , $$
    and $\partial$ is extended to all of $C_n(X)$ by linearity.

\enote

Geometrically, we view a simplicial set $X$ as a collection of (possibly degenerate) $n$-simplices $X_n$ for all $n<\omega$.
If we attach the geometric $n$-simplex $\abs{\Delta^n}$ to each $n$-simplex $X_n$ we obtain $X_n\times\abs{\Delta^n}$.
We give each $X_n$ the discrete topology, so this becomes a topological space.
The geometric realization of $X$ as a whole should be the union of all these simplices, but this is not enough.
For one, degenerate simplices should not be counted; and we want to glue the faces of simplices together.

So we define the geometric realization as follows.

\bdefn

    For a simplicial set $X$, we define its {\emphcolor geometric realization} to be the space
    $$ \abs X = \coprod_{n<\omega}X_n\times\abs{\Delta^n}/{\sim} $$
    where each $X_n$ has the discrete topology and $\sim$ is the equivalence generated by the relations
    $$ (x,\delta_ip) \sim (d_ix,p),\qquad (x,\sigma_ip) \sim (s_ix,p) . $$

    Here $\delta_i\colon\abs{\Delta^{n-1}}\to\abs{\Delta^n}$ and $\sigma_i\colon\abs{\Delta^n}\to\abs{\Delta^{n-1}}$ are the face and degeneracy maps respectively
    of the standard simplices.

\edefn

For $x\in X_n$, we denote by $(x,\abs{\Delta^n})$ the image of $\set x\times\abs{\Delta^n}$ in $\abs X$.

\blemm

    If $X$ is a simplicial set, then every degenerate simplex of $X$ is the degeneracy of a unique non-degenerate simplex.
    That is, if $x\in X$ is degenerate then there is a unique non-degenerate $y\in X$ such that $x=s_{i_1}\cdots s_{i_k}y$ for some $i_j$s.

\elemm

\bproof

    Clearly if $x$ is degenerate there must exist {\it some\/} $y$ which satisfies this.
    Now suppose $x=Sy=S'y'$ for $y,y'$ non-degenerate and $S,S'$ a composition of degeneracy maps.
    Suppose $S=s_{i_1}\cdots s_{i_k}$, then define $D=d_{i_k}\cdots d_{i_1}$.
    This is a left inverse for $S$, and thus $y=DS'y'$.

    Degeneracy and face maps commute according to the identities laid out in \gotopasteqn{simpid}.
    Thus $y=DS'y'=\tilde S\tilde Dy'$ for some composition of degeneracy maps and face maps.
    Since $y$ is by assumption non-degenerate, $\tilde S$ must be empty so that $y=\tilde Dy'$.

    By symmetry, $y'=\bar Dy$ for some composition of face maps as well.
    Since face maps lower the dimension, this can only happen if $\bar D,\tilde D$ are empty.
    Thus $y=y'$.
    \qed

\eproof

Similarly we have the following result.

\blemm

    Every point $v\in\abs{\Delta^n}$ can be expressed uniquely as $v=\delta_{i_1}\cdots\delta_{i_k}u$ where $u$ is an interior point (all of its coordinates are positive).

\elemm

Call a point $(x,v)\in\abs X$ {\emphcolor non-degenerate} if $x$ is non-degenerate and $v$ is interior.

\blemm

    For a simplicial set $X$, every point in its geometric realization $(x,v)\in\abs X$ is equivalent to a unique non-degenerate point of $\abs X$.

\elemm

\bproof

    Define $\lambda\colon\abs X\to\abs X$ as follows.
    For $x\in X_n$ take $i_1,\dots,i_k$ and $y$ non-degenerate such that $x=s_{i_1}\cdots s_{i_k}y$.
    Then define
    $$ \lambda(x,v) = (y,\sigma_{i_k}\cdots\sigma_{i_1}v) . $$
    Similarly define $\mu\colon\abs X\to\abs X$, which for $(x,v)$ maps
    $$ \mu(x,v) = (d_{i_k}\cdots d_{i_1}x,u) $$
    where $v=\delta_{i_1}\cdots\delta_{i_k}u$.

    Clearly $\lambda,\mu$ both map points to equivalent points.
    Then $\lambda\circ\mu$ maps a point to an equivalent non-degenerate point.

    Uniqueness is because it can be verified that if $p\sim p'$ then $\lambda\mu(p)=\lambda\mu(p')$.
    \qed

\eproof

An important fact is the following.

\bthrm

    The geometric realization of a simplicial set $X$ is a CW complex.

\ethrm

\bproof

    Take the $n$-cells of $\abs X$ to be the images of the non-degenerate $n$-simplices of $X$.
    That is, simplices $(x,\abs{\Delta^n})$ for $x\in X_n$ non-degenerate.
    By the previous lemma, these cover $\abs X$, and it can be verified that these make $\abs X$ into a CW complex.
    \qed

\eproof

\bthrm

    There is an adjunction $\abs\bul\colon\sSet\tof\Top\colon\S$ between the geometric realization and the singular set functors.

\ethrm

\bproof

    For a simplicial set $X$ and a topological space $Y$, we define maps
    $$ \Phi\colon\hom_\Top(\abs X,Y) \tof \hom_\sSet(X,\S(Y))\colon\Psi $$
    as follows.
    For a continuous $f\colon\abs X\to Y$, restricting $f$ to a non-degenerate simplex $(x,\abs{\Delta^n})$ restricts to a continuous map $\Phi(f)_n\colon\abs{\Delta^n}\to Y$.

    Conversely for a map $g\colon X\to\S(Y)$, for each $n$-simplex $x\in X$ we have a continuous function $\rho_x\colon\abs{\Delta^n}\to Y$.
    So define $\Psi(g)\colon\abs X\to Y$ to act on the simplex $(x,\abs{\Delta^n})\in\abs X$ by applying $\rho_x$ to $\abs{\Delta^n}$.
    \qed

\eproof

